\documentclass{article}
\usepackage{amsmath,amssymb}
\usepackage[margin=1in]{geometry}
\begin{document}

\section*{FP purchase subsidy with endogenous pricing}

\[
\boxed{\max\bigl\{s_{FP}^{*}(r_P):b_{FP}(r_P)\leq b\bigr\}.}
\]
The sponsor chooses the per-unit purchase subsidy $r_P$; $b$ is her fixed available budget. The firm subsequently chooses its distribution quantity and price to maximize expected profit, before demand is realized. The subscripts $F$ and $P$ designate the for-profit firm and purchase subsidy, respectively. The sponsor's objective and budget constraint above are the unnumbered formulation in Section 5.1 (physical PDF page 20).

\subsection*{Firm response and sponsor accounting}

The purchase-subsidy replacements for the newsvendor quantity and price responses in Section 3.1 are
\begin{align*}
p_{FP}^{*}(r_P)
&=\operatorname*{argmax}_{p\in[\underline p,\bar p]}
\left\{p\int_{\frac{c-r_P}{p}}^{1}
\bar F_D^{inv}(\theta,p)\,d\theta\right\},\\
q_{FP}^{*}(r_P)
&=\bar F_D^{inv}\!\left(
\frac{c-r_P}{p_{FP}^{*}(r_P)},p_{FP}^{*}(r_P)
\right),\\
s_{FP}^{*}(r_P)
&=\hat s\!\left(q_{FP}^{*}(r_P),p_{FP}^{*}(r_P)\right),\\
b_{FP}(r_P)&=r_Pq_{FP}^{*}(r_P).
\end{align*}
These are the operative responses and accounting definitions on physical PDF page 19, rather than the superseded no-subsidy responses (4)--(6). As specified in footnote 5 on physical PDF page 12, throughout the paper, if multiple optimal solutions exist, the smallest one is chosen.

The fixed-price quantity response used in this construction is stated in Section 4.1 (physical PDF page 14):
\[
q_{FP}(r_P,p)=\bar F_D^{inv}\!\left(\frac{c-r_P}{p},p\right).
\]
Here $\bar F_D^{inv}(\cdot,p)$ is the inverse of $\bar F_D(\cdot,p)$, as defined immediately after (4). The paper leaves this inverse notation abstract.

\subsection*{Demand, consumption, profit, and variable roles}

The fixed demand family consists of nonnegative random variables $D(p)$ with distribution $F_D(\cdot,p)$ and mean $\mu_D(p)$, indexed by the firm's admissible price:
\begin{gather*}
p\in[\underline p,\bar p],\qquad
\mu_D(p)=\mathbb E[D(p)],\\
\bar F_D(\cdot,p)=1-F_D(\cdot,p),\\
\bar d(p)=\sup\{D(p)\}<\infty,
\qquad p\in[\underline p,\bar p].
\end{gather*}
Thus $\underline p$ and $\bar p$ are fixed price endpoints, whereas $\bar d(p)$ is the upper endpoint of demand support, not a printed upper bound on the firm's quantity decision. The marginal cost $c$ and sponsor budget $b$ are fixed data. The firm's distribution quantity is $q$, its price is $p$, and $d$ denotes realized demand. The subsidy and distribution quantity are nonnegative quantities; no integrality restriction is stated:
\[
r_P\geq0,\qquad q\geq0,\qquad p\in[\underline p,\bar p].
\]

The purchase subsidy pays $r_P$ dollars for each unit brought to the market. The realized profit is
\begin{equation}
\pi_P(q,p,d,r_P)=p(q\wedge d)-cq+r_Pq,
\tag{1}
\end{equation}
where
\[
x\wedge y=\min\{x,y\}.
\]
The for-profit firm's objective is to maximize
\[
\mathbb E[\pi_P(q,p,D(p),r_P)].
\]
There are no not-for-profit break-even or loss-probability constraints in this formulation.

The expected-sales function defined in Section 4 (physical PDF page 13) is
\[
\hat s(q,p)=\int_0^q\bar F_D(x,p)\,dx
           =\mathbb E[q\wedge D(p)].
\]
The proof of Proposition 3 (physical PDF page 27) writes the corresponding expected-profit function as
\[
\phi(p,q,r_P)=p\int_0^q\bar F_D(x,p)\,dx-(c-r_P)q.
\]
This is the firm's profit objective underlying the response formulas, not the sponsor's consumption objective.

\subsection*{Scope of the monotonicity premise}

The sponsor formulation is stated without a hazard-order constraint. Proposition 3 provides a conditional monotonicity result for the same responses. Its printed premise is that
\[
\{D(p),\ \underline p\leq p\leq\bar p\}
\quad\text{is decreasing in the hazard rate order},
\]
with the printed parenthetical explanation
\[
-\frac{d\log\bar F_D(x,p)}{dx}
\quad\text{is decreasing in $p$ for any $x$}.
\]
Under that premise, the proposition states, for
\[
0\leq r_P\leq c,
\]
that $q_{FP}^{*}(r_P)$, $s_{FP}^{*}(r_P)$, and $b_{FP}(r_P)$ are increasing in $r_P$, while $p_{FP}^{*}(r_P)$ is decreasing in $r_P$. This range is the printed scope of Proposition 3; the sponsor's displayed maximization does not itself print an upper bound $r_P\leq c$.

\paragraph{Printed hazard-order inconsistency.}
The parenthetical hazard-rate direction above is retained as printed. In the proof of Proposition 3 the paper instead uses the implication
\[
\frac{\bar F_D(x,p)}{\bar F_D(q,p)}
\quad\text{is increasing in $p$ for any $x<q$}.
\]
This implication has the opposite hazard-rate direction from the printed parenthetical. Neither reading is silently substituted for the other, and neither is added as an unconditional feasibility restriction on the sponsor's problem.

\subsection*{Full-subsidy boundary and the printed extension}

Section 5.3 and the proof of Proposition 6 identify the full-cost purchase subsidy $r_P=c$. With
\[
p_\mu\equiv\operatorname*{argmax}
\{p\mu_D(p):\underline p\leq p\leq\bar p\},
\]
they state the boundary consumption and the minimum budget needed to induce it as
\[
s_{FP}^{*}(c)=\mu_D(p_\mu),\qquad
b_{FP}(c)=c\bar d(p_\mu).
\]

\paragraph{Printed extension beyond full cost.}
The proof of Proposition 6 (physical PDF page 29) states that, if $r_P\geq c$, the firm's objective is increasing in $q$ and then prints
\[
q_{FP}^{*}(r_P)=\bar d(p),\qquad
p_{FP}^{*}(r_P)=p_\mu,
\]
asserting that the decisions become independent of $r_P$ once $r_P\geq c$. The $p$ in the first expression is retained as printed. For $r_P>c$, this assertion conflicts with (1) when no quantity cap is imposed: after demand is fully covered, profit still increases with $q$. The paper does not state a decision constraint $q\leq\bar d(p)$ in its model definition. Accordingly, no such cap, clipping of the inverse argument, or repaired response for $r_P>c$ is introduced here. The inverse responses are retained as printed; the paper also supplies no explicit extension of $\bar F_D^{inv}$ outside its probability argument range or additional price restriction to resolve such cases.

\end{document}
